\documentclass[11pt,a4paper]{beamer}
\usetheme{Madrid}
\usepackage[utf8]{inputenc}
\usepackage{amsmath}
\usepackage{amsfonts}
\usepackage{amssymb}
\usepackage{graphicx}
\usepackage{xcolor}
\usepackage{tikz}
\usepackage{booktabs}
\usepackage{array}
\usepackage{colortbl}
\usepackage{multicol}

% ============================================================
% EXTREMELY SMALL FONT - MAXIMUM CONTENT PER SLIDE
% ============================================================
\usefonttheme{professionalfonts}

\setbeamerfont{normal text}{size=\scriptsize}
\setbeamerfont{itemize/enumerate body}{size=\scriptsize}
\setbeamerfont{itemize/enumerate subbody}{size=\tiny}
\setbeamerfont{block body}{size=\scriptsize}
\setbeamerfont{block title}{size=\scriptsize}
\setbeamerfont{frametitle}{size=\small}
\setbeamerfont{title}{size=\normalsize}
\setbeamerfont{subtitle}{size=\small}
\setbeamerfont{author}{size=\tiny}
\setbeamerfont{date}{size=\tiny}

\renewcommand{\scriptsize}{\fontsize{6pt}{7pt}\selectfont}
\renewcommand{\footnotesize}{\fontsize{6.5pt}{8pt}\selectfont}
\renewcommand{\small}{\fontsize{7.5pt}{9pt}\selectfont}
\renewcommand{\normalsize}{\fontsize{8.5pt}{10.5pt}\selectfont}

\newcommand{\redflag}[1]{\textcolor{red}{\textbf{[!] #1}}}
\newcommand{\bluenote}[1]{\textcolor{blue}{\textbf{[i] #1}}}
\newcommand{\greennote}[1]{\textcolor{green!50!black}{\textbf{[CORRECT] #1}}}
\newcommand{\examfocus}[1]{\textcolor{orange!80!black}{\textbf{[FOCUS] #1}}}
\newcommand{\trap}[1]{\textcolor{purple}{\textbf{[TRAP] #1}}}
\newcommand{\correct}[1]{\textcolor{green!50!black}{\textbf{[right] #1}}}
\newcommand{\scenario}[1]{\textcolor{blue!70!black}{\textbf{[SCENARIO] #1}}}
\newcommand{\keyterm}[1]{\textcolor{red!80!black}{\textbf{[KEY] #1}}}

\title[CAMS Domain B - Missing Hard Questions]{CAMS Exam: Domain B Missing Content\\15 Hard Questions Missing from Standard Materials}
\author{Based on 2024-2026 Candidate Feedback}
\date{}

\setbeamercovered{transparent}
\setbeamertemplate{navigation symbols}{}
\setbeamertemplate{frametitle continuation}{}
\setlength{\parskip}{2pt}
\setlength{\itemsep}{1pt}

\begin{document}

% ============================================================
% TITLE SLIDE
% ============================================================
\begin{frame}
\titlepage
\centering
\tiny FATF Rec. 35, Rec. 31, Rec. 34, IO.11, 7AMLD, Wolfsberg CBDDQ, 314(b) Opt-In, and More
\end{frame}

% ============================================================
% SECTION 1: FATF RECOMMENDATION 35 - SANCTIONS (COMPLETELY MISSING)
% ============================================================

\begin{frame}{FATF Recommendation 35 - Sanctions}
\begin{block}{Rec. 35 - Effective, Proportionate, Dissuasive (EPD)}
All sanctions for AML/CFT violations must meet THREE criteria:
\begin{enumerate}
\item \textbf{Effective} - actually deter violations
\item \textbf{Proportionate} - fit the severity of violation
\item \textbf{Dissuasive} - discourage future violations
\end{enumerate}
\end{block}
\textbf{Sanctions Available For:}
\begin{tabular}{p{0.45\textwidth}|p{0.45\textwidth}}
\hline
\rowcolor{lightgray}
\textbf{Natural Persons (Individuals)} & \textbf{Legal Persons (Corporations)} \\
\hline
Imprisonment & Significant fines \\
Criminal fines & License revocation \\
Personal liability & Debarment from government contracts \\
Director disqualification & Cease-and-desist orders \\
\hline
\end{tabular}
\scenario{A regulator issues a \$500 fine to a bank for a \$50 million AML program failure.}
\trap{The fine exists, so sanctions are "present."}
\correct{The fine is NOT dissuasive. \$500 for \$50M failure violates EPD standard.}
\end{frame}

\begin{frame}{Rec. 35 - Sanctions Types and Application}
\begin{block}{Three Categories of Sanctions}
\begin{enumerate}
\item \textbf{Criminal Sanctions}: Imprisonment, criminal fines (for willful violations)
\item \textbf{Civil Sanctions}: Civil penalties, damages (for negligent violations)
\item \textbf{Administrative Sanctions}: License revocation, cease-and-desist, debarment
\end{enumerate}
\end{block}
\redflag{MUST be available against BOTH individuals AND legal persons (corporations).}
\vspace{0.2cm}
\textbf{What Rec. 35 Requires:}
\begin{itemize}
\item Sanctions for failure to file STRs
\item Sanctions for failure to conduct CDD
\item Sanctions for tipping-off
\item Sanctions for violations by senior management (not just junior staff)
\end{itemize}
\examfocus{Exam asks: "Which type of sanction is available under FATF Rec. 35?" Answer: Criminal, civil, AND administrative.}
\end{frame}

% ============================================================
% SECTION 2: FATF RECOMMENDATION 31 - INVESTIGATIVE POWERS
% ============================================================

\begin{frame}{FATF Recommendation 31 - Investigative Powers}
\begin{block}{Rec. 31 - Powers of Competent Authorities}
Law enforcement and competent authorities must have comprehensive investigative powers.
\end{block}
\textbf{Mandatory Powers:}
\begin{itemize}
\item Compel production of documents and records (including financial records)
\item Conduct searches of persons and premises
\item Take witness statements under oath
\item Obtain financial information from FIs directly (no court order required)
\item Use \textbf{Special Investigative Techniques (SITs)}
\end{itemize}
\textbf{Special Investigative Techniques (SITs):}
\begin{itemize}
\item Undercover operations
\item Electronic surveillance (wiretapping, monitoring)
\item Controlled deliveries (allowing illicit funds to move to identify recipients)
\item Financial investigations parallel to predicate crime investigations
\end{itemize}
\redflag{All SITs must be available for BOTH ML and TF investigations.}
\end{frame}

\begin{frame}{Rec. 31 - Controlled Delivery Exam Scenario}
\scenario{Terrorist Financing Controlled Delivery}
Law enforcement identifies a cash courier transporting \$150,000 from Country A to Country B, suspected of funding a terrorist cell. Rather than arresting the courier and seizing the cash, authorities allow the transaction to proceed while monitoring the funds. The courier delivers the cash to a religious charity in Country B. Law enforcement then arrests both the courier and the charity's director, and seizes records identifying other donors.
\vspace{0.2cm}
\textbf{Q: What investigative technique is being used?}
\begin{itemize}
\item A) Electronic surveillance
\item B) Controlled delivery
\item C) Undercover operation
\item D) Asset forfeiture
\end{itemize}
\trap{This is regular law enforcement; no special technique needed.}
\correct{Correct Answer: B (Controlled delivery). The funds were allowed to reach their intended destination to identify the full network.}
\end{frame}

\begin{frame}{Rec. 31 - Expeditious Access to Records}
\begin{block}{Critical Requirement - NO Unreasonable Delays}
Law enforcement must be able to obtain financial records \textbf{expeditiously} - within days, not months.
\end{block}
\scenario{Prosecutor Unable to Obtain Bank Records}
A prosecutor is investigating a \$20 million corruption case. Bank records can only be obtained through a court order that takes 6-9 months. When records finally arrive, the suspect has already moved the funds through shell companies in three jurisdictions.
\vspace{0.2cm}
\textbf{Q: Which FATF Recommendation is violated?}
\trap{Recommendation 30 (Law Enforcement Powers) - but that's partially correct.}
\correct{Correct Answer: Recommendation 31. The inability to obtain records expeditiously violates the requirement for effective investigative powers. Rec. 30 covers having an FIU and LE structure; Rec. 31 covers specific powers including timely access.}
\end{frame}

% ============================================================
% SECTION 3: FATF RECOMMENDATION 34 - GUIDANCE AND FEEDBACK
% ============================================================

\begin{frame}{FATF Recommendation 34 - Guidance and Feedback}
\begin{block}{Rec. 34 - NOT Optional}
Competent authorities MUST provide guidance and feedback to FIs and DNFBPs.
\end{block}
\textbf{What Must Be Provided:}
\begin{itemize}
\item Red flag indicators for detecting suspicious activity
\item Typologies and emerging ML/TF trends
\item Feedback on STR quality and usefulness
\item Best practices and case studies
\item Outreach and awareness programs
\end{itemize}
\scenario{FIU Disseminates STRs But Never Provides Feedback}
A country's FIU receives 85,000 STRs annually, disseminates 12,000 to law enforcement, but never provides any feedback to reporting institutions. Banks receive no information on whether their STRs were useful, no typology guidance, and no red flag updates.
\vspace{0.2cm}
\textbf{Q: Which FATF Recommendation is violated?}
\trap{Rec. 20 (STR filing) - but banks ARE filing STRs.}
\correct{Correct Answer: Recommendation 34. Guidance and feedback are REQUIRED, not optional best practices.}
\end{frame}

% ============================================================
% SECTION 4: FATF RECOMMENDATION 7 - PROLIFERATION FINANCING (EXPANDED)
% ============================================================

\begin{frame}{FATF Recommendation 7 - Proliferation Financing}
\begin{block}{Rec. 7 - WMD Proliferation Financing}
Countries must implement targeted financial sanctions for proliferation financing under UNSCR 1540, 1718, 1737, 2231, and 2270.
\end{block}
\textbf{Proliferation Financing Red Flags:}
\begin{itemize}
\item Dual-use goods (precision bearings, maraging steel, radiation-hardened chips, vacuum pumps) to sanctioned countries
\item End-use statements inconsistent with goods ("agricultural research" for missile guidance components)
\item Unusual shipping routes or free trade zone transshipment
\item Customer has no legitimate business need for the goods
\item Payment routing through multiple third-party intermediaries
\item Shell company involvement on buyer side
\end{itemize}
\redflag{Dual-use goods + sanctioned jurisdiction + inconsistent end-use = proliferation financing red flag requiring SAR.}
\end{frame}

\begin{frame}{Rec. 7 - Proliferation Financing Exam Scenario}
\scenario{Suspicious Industrial Equipment Export}
A small trading company with no prior industrial history orders \$3.2 million in precision ball bearings rated for 15,000+ RPM. End-use declared as "food processing equipment" for a country under UNSC sanctions for WMD programs. The buyer pays through a shell company in a free trade zone. Shipping documents show routing through three intermediary countries with no apparent economic purpose.
\vspace{0.2cm}
\textbf{What the Exam Tests:} Recognizing proliferation financing red flags.
\vspace{0.2cm}
\textbf{Analysis:}
\begin{enumerate}
\item Precision bearings at 15,000+ RPM = potential missile or centrifuge component
\item "Food processing" end-use is inconsistent (food processing uses lower-grade bearings)
\item Sanctions jurisdiction + FTZ routing + shell company = layering
\item Proliferation financing SAR required
\end{enumerate}
\trap{This is standard trade-based money laundering.}
\correct{Correct Answer: Proliferation financing under UNSCR 1540 and FATF Rec. 7. Dual-use goods + sanctions evasion + inconsistent end-use = PF red flag.}
\end{frame}

% ============================================================
% SECTION 5: FATF IMMEDIATE OUTCOME 11 (IO.11) - MISSING FROM DOMAIN B
% ============================================================

\begin{frame}{FATF Immediate Outcome 11 (IO.11)}
\begin{block}{IO.11 - Counter-Proliferation Financing Measures}
Measures the effectiveness of a country's ability to prevent and stop WMD proliferation financing.
\end{block}
\textbf{What IO.11 Assesses:}
\begin{itemize}
\item Implementation of UNSCR 1540 (WMD non-proliferation)
\item Targeted financial sanctions for designated proliferators
\item Detection of dual-use goods exports to sanctioned programs
\item Coordination between customs, intelligence, and financial authorities
\item Number of proliferation financing investigations and prosecutions
\end{itemize}
\scenario{Country Has Strong ML/TF Controls But Weak PF Controls}
Country Z has excellent ML/TF controls (High on IO.1-IO.10) but has never identified a proliferation financing case. Dual-use goods flow freely through its free trade zones to sanctioned countries. Customs and financial authorities do not coordinate on proliferation risks.
\vspace{0.2cm}
\textbf{Q: Which FATF Immediate Outcome will be rated Low?}
\trap{IO.7 (ML prosecutions) - but ML prosecutions are fine.}
\correct{Correct Answer: IO.11 (Counter-Proliferation Financing). Weak PF controls result in a separate low rating even if ML/TF controls are strong.}
\end{frame}

% ============================================================
% SECTION 6: FATF RECOMMENDATION 23 - TCSP ACTIVITIES (MISSING DETAIL)
% ============================================================

\begin{frame}{FATF Recommendation 23 - Trust and Company Service Providers (TCSPs)}
\begin{block}{Rec. 23 - DNFBP Obligations for TCSPs}
TCSPs must conduct CDD when performing specified activities.
\end{block}
\textbf{Six TCSP Activities Requiring CDD:}
\begin{enumerate}
\item Acting as a \textbf{formation agent} for legal persons
\item Acting as (or arranging for) a \textbf{director or secretary} of a company
\item Providing a \textbf{registered office} address
\item Acting as (or arranging for) a \textbf{nominee shareholder} for another person
\item Acting as (or arranging for) a \textbf{trustee} of an express trust
\item Acting as (or arranging for) a \textbf{nominee director} for another person
\end{enumerate}
\scenario{Law Firm Forms Shell Companies for Anonymous Client}
A law firm forms 15 BVI companies for a client the firm has never met. The firm provides nominee directors and registered office addresses. The client refuses to provide identification or disclose the beneficial owner.
\vspace{0.2cm}
\textbf{Q: What is the law firm's obligation under FATF Rec. 23?}
\trap{Attorney-client privilege exempts the firm from CDD.}
\correct{Correct Answer: The firm is acting as a TCSP (formation agent + nominee director + registered office). CDD is required. Attorney-client privilege does NOT apply to company formation activities.}
\end{frame}

% ============================================================
% SECTION 7: EU 7AMLD AND AMLA (CURRENT EVENTS - MISSING)
% ============================================================

\begin{frame}{EU 7AMLD and AMLA (2024-2026 Updates)}
\begin{block}{EU AML Authority (AMLA)}
First centralized EU AML supervisor with \textbf{direct supervisory powers} over highest-risk cross-border financial institutions.
\end{block}
\textbf{Key 7AMLD / AMLA Provisions:}
\begin{itemize}
\item \textbf{AMLA} operational by 2026, headquartered in Frankfurt
\item Direct supervision of \textbf{high-risk cross-border FIs} (banks, VASPs, MSBs)
\item \textbf{EUR 10,000 cash payment limit} across all EU member states (mandatory)
\item \textbf{Beneficial ownership registers} - public access restored (subject to legitimate interest test)
\item \textbf{VASP Travel Rule} enforcement strengthened
\item \textbf{Supervisory powers} for AMLA: on-site inspections, fines, license revocation
\end{itemize}
\redflag{Exam QUESTION: "What is the new EU AML supervisor with direct powers?" Answer: AMLA (EU AML Authority), not EBA.}
\end{frame}

\begin{frame}{EU 7AMLD - EUR 10,000 Cash Limit}
\begin{block}{Mandatory EU-Wide Cash Payment Limit}
All EU member states must prohibit cash payments exceeding EUR 10,000.
\end{block}
\textbf{Key Points:}
\begin{itemize}
\item Applies to all commercial transactions (goods, services, real estate)
\item Member states may set \textbf{LOWER} limits (e.g., EUR 3,000)
\item Violations subject to sanctions
\item Aims to reduce cash-based ML/TF and tax evasion
\item Does NOT apply to private individuals (e.g., selling personal car)
\end{itemize}
\scenario{A real estate transaction in Spain for EUR 450,000. The buyer offers to pay EUR 350,000 in cash and the remainder via bank transfer.}
\vspace{0.2cm}
\textbf{Q: Under 7AMLD, is this transaction permitted?}
\trap{Cash payments are still permitted as long as the transaction is reported.}
\correct{Correct Answer: NO. EUR 350,000 cash exceeds the EUR 10,000 limit. The entire transaction would violate the cash payment prohibition unless restructured entirely through banking channels.}
\end{frame}

% ============================================================
% SECTION 8: WOLFSBERG CBDDQ - THREE SECTIONS (MISSING DETAIL)
% ============================================================

\begin{frame}{Wolfsberg Correspondent Banking Due Diligence Questionnaire (CBDDQ)}
\begin{block}{The Three Critical Sections of the CBDDQ}
The Wolfsberg CBDDQ is organized into three core areas.
\end{block}
\begin{tabular}{p{0.30\textwidth}|p{0.65\textwidth}}
\hline
\rowcolor{lightgray}
\textbf{Section} & \textbf{What It Covers} \\
\hline
\textbf{Section 1: AML/CFT Policies \& Procedures} & Respondent's AML controls, SAR process, independent audit, CDD policies, recordkeeping \\
\hline
\textbf{Section 2: Customer Base Risk} & Types of customers (MSBs, NBFIs, foreign PEPs, shell companies, high-risk jurisdictions) \\
\hline
\textbf{Section 3: Independent Testing} & Audit schedule, scope, findings, remediation status, auditor qualifications \\
\hline
\end{tabular}
\redflag{Most critical risk factor: Respondent servicing unlicensed MSBs + no recent independent audit.}
\trap{The CBDDQ asks about the respondent bank's CEO personal investments or office locations.}
\correct{Correct Answer: The CBDDQ focuses exclusively on AML controls, customer base composition, and independent testing - not executive personal matters.}
\end{frame}

% ============================================================
% SECTION 9: USA PATRIOT ACT SECTION 312 - EDD FOR CORRESPONDENT ACCOUNTS
% ============================================================

\begin{frame}{USA PATRIOT Act Section 312 - EDD for Correspondent Accounts}
\begin{block}{Section 312 - Mandatory Enhanced Due Diligence}
U.S. banks must apply EDD to correspondent accounts of foreign banks.
\end{block}
\textbf{Triggering Conditions for EDD:}
\begin{itemize}
\item Foreign bank operates under an \textbf{offshore license} (Cayman, BVI, Bahamas, Bermuda, Panama)
\item Foreign bank has \textbf{PEP customers} as a significant customer base
\item Foreign bank located in \textbf{Grey List or Black List} jurisdiction
\item Foreign bank permits \textbf{payable-through accounts} for non-bank customers
\item Foreign bank is NOT a subsidiary of a U.S. bank
\end{itemize}
\textbf{Required EDD Elements (Section 312):}
\begin{enumerate}
\item Determine if account is a payable-through account
\item Identify all owners of the foreign bank
\item Conduct enhanced scrutiny of transactions
\item Determine if the foreign bank has AML controls for its own customers
\end{enumerate}
\redflag{Section 312 EDD is MANDATORY - not risk-based discretion.}
\end{frame}

% ============================================================
% SECTION 10: FINCEN 314(b) - OPT-IN REQUIREMENT (MISSING DETAIL)
% ============================================================

\begin{frame}{USA PATRIOT Act Section 314(b) - Voluntary Information Sharing}
\begin{block}{314(b) - Critical Opt-In Requirement}
Financial institutions must file an \textbf{opt-in notice} with FinCEN before sharing information.
\end{block}
\textbf{314(b) Requirements:}
\begin{itemize}
\item File \textbf{opt-in notice} with FinCEN (not automatic)
\item Maintain \textbf{confidentiality} of shared information
\item Sharing limited to \textbf{suspected ML/TF} (cannot share for other purposes)
\item May share SAR information (exception to tipping-off prohibition)
\item Receive \textbf{safe harbor} from liability
\end{itemize}
\textbf{Eligible Institutions:}
\begin{itemize}
\item Banks, credit unions, savings associations
\item Broker-dealers, futures commission merchants
\item MSBs, casinos, insurance companies
\item \textbf{NOT available to:} individuals, most DNFBPs (lawyers, accountants, real estate agents generally excluded)
\end{itemize}
\trap{Any financial institution can share information under 314(b) without prior registration.}
\correct{Correct Answer: FI must FILE OPT-IN NOTICE with FinCEN before sharing. Retroactive sharing is not permitted.}
\end{frame}

% ============================================================
% SECTION 11: EU 5AMLD - PREPAID CARD THRESHOLD (MISSING NUMBER)
% ============================================================

\begin{frame}{EU 5AMLD - Prepaid Card Threshold}
\begin{block}{Lowered Threshold for Anonymous Prepaid Cards}
5AMLD lowered the anonymous prepaid card threshold from EUR 250 to EUR 150.
\end{block}
\textbf{Key Provisions:}
\begin{itemize}
\item 4AMLD threshold: EUR 250
\item 5AMLD threshold: \textbf{EUR 150} (effective 2020)
\item Prepaid cards with value exceeding EUR 150 must be linked to an identified customer account
\item Applies to all anonymous prepaid cards (including digital/virtual cards)
\item Applies to single-use and reloadable cards
\end{itemize}
\scenario{Customer Purchases Multiple EUR 100 Prepaid Cards}
A customer purchases 40 prepaid debit cards, each loaded with EUR 100 (total EUR 4,000). Each individual card is below the EUR 150 threshold.
\vspace{0.2cm}
\textbf{Q: Under 5AMLD, is this permitted without customer identification?}
\trap{Each card is below EUR 150, so no identification required.}
\correct{Correct Answer: The 5AMLD threshold applies PER CARD. EUR 100 per card is below the threshold, so no identification required for individual cards. However, structural purchases may trigger other AML concerns (structuring).}
\redflag{Exam QUESTION: "Under 5AMLD, what is the maximum value of an anonymous prepaid card?" Answer: EUR 150.}
\end{frame}

% ============================================================
% SECTION 12: UNSCR 1373 vs. UNSCR 1267/1988 - KEY DISTINCTION
% ============================================================

\begin{frame}{UNSCR 1373 vs. UNSCR 1267/1988 - Critical Difference}
\begin{block}{Two Different Types of Terrorist Financing Sanctions}
The exam tests the distinction between UNSCR 1373 and UNSCR 1267/1988.
\end{block}
\begin{tabular}{p{0.45\textwidth}|p{0.45\textwidth}}
\hline
\rowcolor{lightgray}
\textbf{UNSCR 1267/1988} & \textbf{UNSCR 1373} \\
\hline
Specific designated persons (Al-Qaida, Taliban, ISIS) & General obligation to criminalize TF \\
UN Security Council maintains the list & Countries create their OWN designation lists \\
\textbf{Mandatory freezing} for all member states & Countries have discretion to designate \\
Assets frozen immediately upon designation & Countries must freeze assets of domestically designated persons \\
No country discretion to delist & Countries can establish their own delisting procedures \\
\hline
\end{tabular}
\scenario{Bank Customer Appears on No Sanctions List}
A bank's customer is not listed on any UN sanctions list. However, the country's domestic counter-terrorism authority has designated the individual as a terrorist financier.
\vspace{0.2cm}
\textbf{Q: Must the bank freeze the customer's assets?}
\trap{No UN listing means no freezing obligation.}
\correct{Correct Answer: YES - under UNSCR 1373, countries must freeze assets of DOMESTICALLY designated persons. The bank must comply with its national law implementing UNSCR 1373.}
\end{frame}

% ============================================================
% SECTION 13: BSA FIVE PILLARS - CDD AS FIFTH PILLAR (MISSING)
% ============================================================

\begin{frame}{BSA Five Pillars - The 2016 CDD Rule}
\begin{block}{Original Four Pillars (pre-2016)}
\begin{enumerate}
\item Internal controls reasonably designed to assure BSA compliance
\item Designated compliance officer (BSA Officer/MLRO)
\item Ongoing employee training (initial and periodic)
\item Independent testing (audit by independent party)
\end{enumerate}
\end{block}
\begin{block}{Fifth Pillar Added in 2016 (FinCEN CDD Rule)}
\textbf{\textcolor{red}{5. Customer Due Diligence (including Beneficial Ownership Identification)}}
\end{block}
\textbf{What the CDD Rule Requires:}
\begin{itemize}
\item Identify and verify beneficial owners of legal entity customers
\item Two-prong test: Ownership (25\%+) AND Control (at least one individual)
\item Applies to all legal entity customers (LLCs, corporations, partnerships, trusts)
\end{itemize}
\redflag{Exam QUESTION: "What regulation added beneficial ownership to the BSA five pillars?"}
\correct{Correct Answer: The 2016 CDD Rule (FinCEN). Beneficial ownership identification is now the FIFTH pillar, not optional.}
\end{frame}

% ============================================================
% SECTION 14: EGMONT GROUP MEMBERSHIP CRITERIA (MISSING)
% ============================================================

\begin{frame}{Egmont Group - Membership Criteria}
\begin{block}{Five Requirements for Egmont Membership}
An FIU must meet ALL of the following criteria.
\end{block}
\textbf{Egmont Membership Requirements:}
\begin{enumerate}
\item FIU must be \textbf{operational} (receiving, analyzing, disseminating intelligence)
\item Country must be a member of \textbf{FATF OR an FSRB} (Asia/Pacific Group, MONEYVAL, GAFILAT, etc.)
\item Must sign the \textbf{Egmont Group Charter}
\item Must have \textbf{legal authority} to exchange information with foreign FIUs
\item Must maintain \textbf{confidentiality} of shared information
\end{enumerate}
\scenario{A country has an operational FIU but is not a member of FATF or any FSRB. The FIU applies for Egmont Group membership.}
\vspace{0.2cm}
\textbf{Q: Will the Egmont Group accept this FIU?}
\trap{Any operational FIU can join Egmont regardless of FATF status.}
\correct{Correct Answer: NO. The country must be a member of FATF OR an FSRB as a prerequisite for Egmont membership.}
\end{frame}

% ============================================================
% SECTION 15: BASEL COMMITTEE - THREE LINES OF DEFENSE (EXPANDED)
% ============================================================

\begin{frame}{Basel Committee - Three Lines of Defense Model}
\begin{block}{Integrated AML/CFT Governance}
The Basel Committee expects AML/CFT to be integrated across all three lines of defense.
\end{block}
\begin{tabular}{p{0.25\textwidth}|p{0.70\textwidth}}
\hline
\rowcolor{lightgray}
\textbf{Line} & \textbf{Responsibility} \\
\hline
\textbf{First Line: Business Operations} & Front-line staff, relationship managers, tellers. Own customer relationships, identify anomalies, collect CDD. \\
\hline
\textbf{Second Line: Compliance} & Compliance Department, BSA Officer, MLRO. Oversight, guidance, monitoring, reporting, SAR filing. \\
\hline
\textbf{Third Line: Internal Audit} & Internal Audit Department. Independent assurance to the Board. Cannot design controls they later test. \\
\hline
\end{tabular}
\scenario{Internal Audit designs the AML monitoring system and also tests it annually.}
\trap{This is efficient - same team knows the system well.}
\correct{Correct Answer: This violates independence requirements. The third line (Audit) cannot design the controls it later tests.}
\redflag{Independence violation: Audit cannot design controls they later test.}
\end{frame}

% ============================================================
% SECTION 16: GDPR ARTICLES 23, 48, 49 (MISSING FROM VARIATION_V2)
% ============================================================

\begin{frame}{GDPR Articles for AML Compliance - Beyond Article 17}
\begin{block}{Article 23 - Restrictions on Data Subject Rights}
EU member states may restrict data subject rights (access, erasure, portability) when necessary for AML/CFT purposes.
\end{block}
\begin{block}{Article 48 - Cross-Border Data Transfers for Court Judgments}
Requires that non-EU court judgments ordering data transfers be recognized only through an MLAT or international agreement.
\end{block}
\begin{block}{Article 49(1)(d) - Derogation for Important Public Interest}
Allows cross-border data transfers when "necessary for important public interest" including AML/CFT and law enforcement cooperation.
\end{block}
\scenario{EU Bank Receives U.S. Grand Jury Subpoena}
A U.S. grand jury subpoenas an EU bank for customer records related to a money laundering investigation. No MLAT exists between the U.S. and the EU country.
\vspace{0.2cm}
\textbf{Q: Under GDPR, can the bank comply?}
\trap{GDPR prohibits all cross-border transfers without an MLAT.}
\correct{Correct Answer: Article 49(1)(d) allows transfer if "necessary for important public interest" including AML/CFT. The bank must document the legal basis but can comply.}
\end{frame}

% ============================================================
% SECTION 17: OFAC 50% RULE - MULTI-TIER CALCULATION (REFRESHER)
% ============================================================

\begin{frame}{OFAC 50\% Rule - Multi-Tier Ownership Calculation}
\begin{block}{How to Calculate Through Multiple Tiers}
Ownership is calculated by multiplying through each tier. Only blocked entities transmit blocked status downstream.
\end{block}
\textbf{Example Structure:}
\begin{itemize}
\item Individual X (SDN-listed) owns 55\% of Company A
\item Company A owns 40\% of Company B
\item Company B owns 60\% of Company C
\end{itemize}
\textbf{Analysis:}
\begin{enumerate}
\item Company A is BLOCKED (55\% owned by SDN - exceeds 50\%)
\item Company B is NOT BLOCKED (40\% ownership from A - below 50\%)
\item Company C is NOT BLOCKED (Company B is not blocked, so its 60\% ownership does NOT transmit blocked status)
\end{enumerate}
\redflag{If an entity in the chain is not blocked, it does NOT transmit blocked status downstream, even if ownership exceeds 50\%.}
\trap{Company C is blocked because Company B owns 60\% of it.}
\correct{Correct Answer: Company C is NOT blocked. Only blocked entities transmit blocked status. Since Company B is not blocked, it cannot transmit blocked status to Company C.}
\end{frame}

% ============================================================
% CONCLUSION SLIDES
% ============================================================

\begin{frame}{Domain B Missing Content - Summary}
\begin{multicols}{2}
\begin{enumerate}
\item \textbf{Rec. 35} - EPD Sanctions (criminal, civil, administrative)
\item \textbf{Rec. 31} - Investigative Powers (SITs: controlled delivery, surveillance)
\item \textbf{Rec. 34} - Guidance \& Feedback (NOT optional)
\item \textbf{Rec. 7} - Proliferation Financing (dual-use goods red flags)
\item \textbf{IO.11} - Counter-Proliferation (separate from ML/TF)
\item \textbf{Rec. 23} - TCSP Activities (6 specific activities)
\item \textbf{7AMLD/AMLA} - EUR 10k cash limit, AMLA supervisor
\item \textbf{Wolfsberg CBDDQ} - 3 sections (controls, customers, testing)
\item \textbf{Section 312} - EDD for correspondent accounts
\item \textbf{Section 314(b)} - Opt-in notice required
\item \textbf{5AMLD Prepaid} - EUR 150 threshold
\item \textbf{UNSCR 1373 vs 1267} - Domestic vs UN designation
\item \textbf{BSA Fifth Pillar} - 2016 CDD Rule (BO identification)
\item \textbf{Egmont Membership} - FATF/FSRB membership required
\item \textbf{Basel Three Lines} - Independence requirement
\item \textbf{GDPR Arts 23,48,49} - AML exceptions to data protection
\item \textbf{OFAC 50\% Rule} - Multi-tier calculation
\end{enumerate}
\end{multicols}
\greennote{These 17 concepts are consistently reported as "surprise" questions.}
\end{frame}

\begin{frame}{Critical Red Flags - Domain B Missing Content}
\begin{itemize}
\item \$500 fine for \$50M violation → NOT dissuasive (violates Rec. 35)
\item 6-month delay for bank records → violates Rec. 31 (expeditious access)
\item FIU never provides feedback to banks → violates Rec. 34
\item Dual-use goods + "food processing" end-use → proliferation financing (Rec. 7)
\item Weak PF controls despite strong ML/TF → Low IO.11 rating
\item Law firm forms companies without CDD → TCSP activity (Rec. 23)
\item EUR 350,000 cash for real estate → violates 7AMLD EUR 10k limit
\item Respondent bank servicing unlicensed MSBs → Wolfsberg CBDDQ highest risk
\item Foreign offshore bank correspondent → Section 312 EDD required
\item 314(b) sharing without opt-in notice → violation
\item Anonymous prepaid card at EUR 200 → violates 5AMLD (limit EUR 150)
\item Country designates TF without UN listing → UNSCR 1373 permits this
\item AML program missing BO identification → missing fifth pillar (2016 CDD Rule)
\item FIU not in FATF or FSRB → cannot join Egmont
\item Audit designs AND tests AML system → violates independence
\item Bank refuses U.S. subpoena citing GDPR → Article 49 provides exception
\end{itemize}
\end{frame}

\begin{frame}{Exam Day Reminders - Domain B Missing Questions}
\begin{itemize}
\item \textbf{Rec. 35}: Sanctions must be Effective, Proportionate, Dissuasive (all three)
\item \textbf{Rec. 31}: Controlled deliveries are permitted special investigative techniques
\item \textbf{Rec. 34}: Guidance and feedback to FIs is MANDATORY, not optional
\item \textbf{Rec. 7 / IO.11}: Proliferation financing is separate from TBML
\item \textbf{TCSPs}: Forming companies, nominee directors, registered office all require CDD
\item \textbf{7AMLD}: EUR 10,000 cash limit across all EU; AMLA is new supervisor
\item \textbf{Wolfsberg CBDDQ}: Three sections = Policies, Customer Base, Independent Testing
\item \textbf{Section 312}: EDD for correspondent accounts of offshore banks is MANDATORY
\item \textbf{Section 314(b)}: Must file OPT-IN NOTICE before sharing
\item \textbf{5AMLD prepaid}: Anonymous cards limited to EUR 150
\item \textbf{UNSCR 1373}: Countries can create their own TF designation lists
\item \textbf{BSA Five Pillars}: 2016 CDD Rule added BO identification as 5th pillar
\item \textbf{Egmont}: Country must be FATF OR FSRB member to join
\item \textbf{Basel Three Lines}: Audit cannot design controls it tests (independence)
\item \textbf{GDPR}: Articles 23, 48, 49 provide AML exceptions to data protection
\item \textbf{OFAC 50\%}: Multiply through tiers; only blocked entities transmit status
\end{itemize}
\vspace{0.3cm}
\centering
\greennote{Good luck - you've covered the hard Domain B questions now!}
\end{frame}

\end{document}